\cotefile{81}
\scanpage{19}
Dém. par récurrence, c'est trivial si $\operatorname{card} I=0,1$ Le cas $\operatorname{card} I=2$ \unsure{a} été vu. Cas général \unsure{avec} $\operatorname{card} I\geq 3$, on choisit $J$ comme dessus. Il suffit de voir que les ens.
\[
\unsure{U|\deleted{=}U}_{J,i}\cap\Sigma_i^+,\ U_{J,i}\cap\Sigma_i^-
\]
ont la forme précédente, il suffit de le voir pour \unsure{l'une|l'un} des deux.

Par hyp. de récurrence, \unsure{on a}
\[
U_{J,i}=X\setminus\coprod_{\alpha\in \deleted{\unsure{J}}K_0}\widehat{\Sigma}_\alpha^+
\]
où
\[
\left\{\begin{aligned}
&\deleted{\Sigma_{\unsure{\alpha|0}}}K_0\subset J\\
&\Sigma_\alpha^+\in\unsure{\pi_0|\Pi_0}(X\setminus E_\alpha)\qquad (\forall\,\alpha\in K_0)\\
&\widehat{\Sigma}_\alpha^+=\Sigma_\alpha^+\cup E_\alpha.
\end{aligned}\right.
\]
\deleted{\unsure{D|O}}\unsure{On a} donc
\[
U=U_{J,i}\cap\Sigma_i^-=X\setminus\left(\coprod_{\alpha\in K_0}\widehat{\Sigma}_\alpha^+\ \cup\ \widehat{\Sigma}_i^+\right)
\]
\deleted{Soit} Au deuxième $=$, dans la \deleted{\unsure{somme|réunion}} réunion $\bigcup_{\alpha\in K_0}\widehat{\Sigma}_\alpha^+$, il suffit de se borner \unsure{à|:} \unsure{sommer} sur les \deleted{$\Sigma$} $\alpha\in K_0$ tels que $\widehat{\Sigma}_\alpha^+\not\subset\widehat{\Sigma}_i^+$ i.e. $\Sigma_\alpha^+\not\subset\Sigma_i^+$ soit $K_1$ cet ens., donc
\[
U=X\setminus\left(\bigcup_{\alpha\in K_1}\widehat{\Sigma}_\alpha^+\ \cup\ \widehat{\Sigma}_i^+\right)
\]
% archive sheet number: 18
